% CONTRACT
%   Establishes: the audit procedure, reproducibly, so another group can run it.
%   Sources: reproduction_manifest (the command column is the ground truth for
%     every procedure described here).
%   May not claim: that the 55 existing tests are held-out evaluation. They are
%     development material (decision record §4, week 1).

\subsection{The claim-check inventory}
We assembled the validation checks the project had identified as unable to fail
and classified each into logical impossibility, low power, implementation error,
or provenance/reporting error. The inventory carries 24 entries against the 21
the project's prose asserted; transcribing the ordinals scattered through the
documentation gives 23, and one further entry was found by running new code
rather than by review. Of the 24, 10 are implementation errors, 6 logical
impossibilities, 6 provenance/reporting errors and 2 low-power; 15 name a
committed input that forces the check to fail, and one is disputed.

\textbf{One worker classified them, and inter-rater agreement is not
established.} The table records a single rater and no second classification. A
rater compared with themselves is a check that cannot fail, so the inventory
does not claim reliability it does not have. The classification is a starting
taxonomy, not a field-wide prevalence estimate.

\subsection{The challenge protocol}
A blinded challenger constructs previously unseen clean and violating cases
against the checks, seals each expectation before running it, and classifies the
consequence as decision-relevant, input-validation or algebraic-identity. A case
counts toward the stop rule only when it was constructed without reading the
ledger, so the existing twelve cases do not count and neither do the project's
own regression tests. The protocol is described here as a design; Results reports
that it has not been executed.

\subsection{Stress calibration}
We hold the Poisson detection model, the cohorts and the nominal level fixed and
vary only the distribution of the per-patient effect: gaussian, skewed,
heavy-tailed, spike-slab, and two boundary regimes where detection is near zero or
near one. Each shape is standardised to unit variance and given mean zero, so the
between-patient scale is the same everywhere and every regime is a true null.
Seven cohorts are kept separate, with six generators, three interval methods,
five fixed seeds and 800 trials per cell; each cell carries a binomial standard
error and a Wilson interval.

\subsection{Bracket reporting}
A calibration returns the smallest bin-median mature-cell count meeting a
criterion, which is a point on a coarse grid. We report instead the interval
between the last failing and first passing bin, and we keep the two criteria
apart: \emph{wide} is coverage alone, \emph{ok} is coverage and discrimination.
Reporting them separately is what makes a crossing that the point rule discards
visible at all.

\subsection{Sensitivity design}
The descriptive adenoma summaries are recomputed across all three weightings,
four granularity rungs, both denominators and all four scale-free statistics,
with patient counts and missingness, and with the estimability gates counted
separately. No unestimable term enters a contrast as a zero.

\subsection{Table resolution}
Every number in this paper is read from a versioned result table by its sidecar
timestamp, not by directory name; where a table has more than one same-date run
the choice is recorded rather than implied. The reproduction index lists the
producing command for each deliverable.
